\documentclass[11pt]{article}
\usepackage[margin=1in]{geometry}
\usepackage{amsmath,amssymb,mathtools}
\usepackage{booktabs}
\usepackage{longtable}
\usepackage{enumitem}
\usepackage{microtype}
\usepackage{xcolor}
\usepackage{caption}
\usepackage{listings}
\usepackage[colorlinks=true,linkcolor=blue!50!black,citecolor=blue!50!black,urlcolor=blue!50!black]{hyperref}
\usepackage{parskip}

\newcommand{\E}{\mathbb{E}}
\newcommand{\ind}[1]{\mathbb{1}_{\{#1\}}}
\newcommand{\sgn}{\operatorname{sgn}}
\DeclareMathOperator{\Var}{Var}
\DeclareMathOperator{\Corr}{Corr}

\definecolor{codebg}{gray}{0.97}
\definecolor{codekw}{rgb}{0.15,0.15,0.55}
\definecolor{codecm}{rgb}{0.35,0.45,0.35}
\lstset{
  language=Python,
  basicstyle=\ttfamily\footnotesize,
  keywordstyle=\color{codekw}\bfseries,
  commentstyle=\color{codecm}\itshape,
  backgroundcolor=\color{codebg},
  frame=single,
  rulecolor=\color{gray!40},
  breaklines=true,
  showstringspaces=false,
  columns=fullflexible,
  captionpos=b,
  xleftmargin=1em,
  framexleftmargin=1em,
}

\title{\textbf{Order-Flow and Imbalance Features}\\[4pt]
\large Formulas, estimation settings and reference implementations\\
from ten source papers, fitted dynamically on data}
\author{Compiled from the source PDFs}
\date{September 2026}

\begin{document}
\maketitle

\begin{abstract}
\noindent This document states how order-flow imbalance (OFI) and related
microstructure features are constructed, exactly as each source paper defines
them, and gives reference implementations that are fitted on rolling windows
with point-in-time inputs only. Every formula is transcribed from the source
PDF, with evident typographical errors in the published versions corrected
here. Where a paper leaves a step genuinely unstated, the document says so and
names the choice the implementer has to make.
Section~\ref{sec:notation} fixes notation; section~\ref{sec:protocol}
gives the shared fitting protocol; sections~\ref{sec:mlofi}--\ref{sec:genimb}
cover the five order-book constructions and the trade-only construction;
sections~\ref{sec:petit}--\ref{sec:bonds} cover feature construction by
clustering and the two prerequisites for applying any of this to corporate
bonds. Section~\ref{sec:map} maps each result back to its paper.
\end{abstract}

\tableofcontents

%======================================================================
\section{Data model, clocks and notation}\label{sec:notation}

A \emph{book event} is any message that changes the visible limit order book: a
limit order arrival, a cancellation or deletion, or an execution. Events are
indexed by $n$ in arrival order with timestamps $\tau_n$, and
\[
  N(t)=\max\{\,n:\tau_n\le t\,\}
\]
is the index of the last event no later than $t$. A \emph{bucket} of length $h$
ending at $t$ is $(t-h,\,t]$, containing events $n=N(t-h)+1,\dots,N(t)$, and
$\Delta N(t)=N(t)-N(t-h)$ counts them.

At level $m$ after event $n$, write $P^{m,b}_{n}$, $P^{m,a}_{n}$ for the bid and
ask prices and $q^{m,b}_{n}$, $q^{m,a}_{n}$ for the queue sizes; level $m=1$ is
the touch. All book quantities are measured \emph{immediately after} the event
is applied. Only populated levels count, so level $m+1$ need not be one tick
beyond level $m$. For trade data without a book, $\varepsilon_t\in\{-1,+1\}$ is
the trade sign and $q_t$ the traded size. The mid-price is
$P_t=(P^{1,b}_t+P^{1,a}_t)/2$.

Three conventions decide whether an implementation is correct:

\begin{enumerate}[leftmargin=2em,itemsep=2pt]
\item \textbf{Bucket alignment.} A contemporaneous regression puts feature and
return on the \emph{same} bucket. A forecasting regression puts the feature on
$(t-h,t]$ and the return on $(t,t+f]$. Mixing them produces a look-ahead result
that looks excellent.
\item \textbf{Sign convention.} Most papers here use $+1$ for buyer-initiated
flow. Gould et al.\ (section~\ref{sec:support}) use the opposite. Check before
combining series from different sources.
\item \textbf{Timestamp units.} Event times and horizons must be compared in one
unit. Casting a timestamp series in a resolution-dependent way while comparing
it against a fixed-unit scalar silently empties every window.
\end{enumerate}

%======================================================================
\section{Fitting protocol}\label{sec:protocol}

All features below share one estimation pattern, kept separate from the feature
definitions so a new feature can be dropped in without touching the estimator.

\paragraph{Rolling windows.} Parameters are estimated on a trailing window of
$W$ observations ending at $t$ and applied only to later observations. The
source papers use short windows: 30 minutes of 10-second buckets (180
observations) contemporaneously, 30 minutes of 1-minute buckets (30
observations) for forecasting, re-estimated every minute.

\paragraph{Regularisation.} Multi-level OFI regressors are correlated across
levels by construction, and the event-type decomposition triples the regressor
count on the same window. Ordinary least squares over-fits in both cases: Xu et
al.\ report out-of-sample error \emph{rising} beyond about five levels under
OLS but not under ridge, and Sitaru et al.\ report out-of-sample $R^2$ below
$-100\%$ under OLS at six levels or more. The papers use ridge with
cross-validated penalty (Xu), LASSO with cross-validated penalty (Cont,
Sitaru, Hu), or principal components.

\paragraph{Evaluation.} In-sample fit is adjusted $R^2$. Out-of-sample fit is
computed on the next disjoint window. For short-horizon \emph{forecasts},
out-of-sample $R^2$ in this literature is routinely negative even where the
economic evaluation is strongly positive, so a trading-rule evaluation is
reported alongside it (sections~\ref{sec:iofi} and~\ref{sec:dofi}).

\begin{lstlisting}[caption={Shared rolling-fit harness. The estimator is injected, so the
same code serves ridge, LASSO and principal-component variants.}]
import numpy as np
import pandas as pd
from sklearn.linear_model import Ridge, LassoCV, RidgeCV
from sklearn.preprocessing import StandardScaler
from sklearn.pipeline import make_pipeline


def rolling_fit_predict(X, y, window, make_model, step=1):
    """Fit on a trailing window, predict forward. Point-in-time by construction.

    Row t is predicted by a model fitted on rows [t - window, t - 1] only.
    """
    X, y = X.sort_index(), y.reindex(X.sort_index().index)
    out = []
    for end in range(window, len(X), step):
        xtr, ytr = X.iloc[end - window:end], y.iloc[end - window:end]
        ok = xtr.notna().all(axis=1) & ytr.notna()
        if ok.sum() < max(10, X.shape[1] + 2):
            continue                      # not enough rows to identify the fit
        model = make_model().fit(xtr[ok], ytr[ok])
        te = slice(end, min(end + step, len(X)))
        out.append(pd.Series(model.predict(X.iloc[te].fillna(0.0)),
                             index=X.index[te], name="y_hat"))
    return pd.concat(out) if out else pd.Series(dtype=float, name="y_hat")


def oos_r2(y_true, y_hat, benchmark=0.0):
    """Out-of-sample R^2 against a constant benchmark.

    NOTE: none of the source papers writes this formula down. This is the
    standard choice (zero benchmark); state it explicitly in any writeup.
    """
    y_true, y_hat = y_true.align(y_hat, join="inner")
    ss_res = float(((y_true - y_hat) ** 2).sum())
    ss_tot = float(((y_true - benchmark) ** 2).sum())
    return 1.0 - ss_res / ss_tot if ss_tot > 0 else np.nan
\end{lstlisting}

%======================================================================
\section{Multi-level order-flow imbalance (MLOFI)}\label{sec:mlofi}

\textbf{Source.} Xu, Gould \& Howison (2020), \emph{Multi-Level Order-Flow
Imbalance in a Limit Order Book}, Market Microstructure and Liquidity 4(3\&4).

\subsection{Definition}

Write $r^m(\tau_n)$ for the total size at the level-$m$ bid and $q^m(\tau_n)$
for the total size at the level-$m$ ask. The bid- and ask-side changes are
\begin{align}
\Delta W^m(\tau_n)&=
\begin{cases}
r^m(\tau_n), & b^m(\tau_n)>b^m(\tau_{n-1}),\\
r^m(\tau_n)-r^m(\tau_{n-1}), & b^m(\tau_n)=b^m(\tau_{n-1}),\\
-r^m(\tau_{n-1}), & b^m(\tau_n)<b^m(\tau_{n-1}),
\end{cases}
\\[6pt]
\Delta V^m(\tau_n)&=
\begin{cases}
-q^m(\tau_{n-1}), & a^m(\tau_n)>a^m(\tau_{n-1}),\\
q^m(\tau_n)-q^m(\tau_{n-1}), & a^m(\tau_n)=a^m(\tau_{n-1}),\\
q^m(\tau_n), & a^m(\tau_n)<a^m(\tau_{n-1}),
\end{cases}
\end{align}
and the level-$m$ contribution and its aggregate are
\begin{equation}
e^m(\tau_n)=\Delta W^m(\tau_n)-\Delta V^m(\tau_n),
\qquad
\mathrm{MLOFI}^m(t_{k-1},t_k)=\sum_{\{n\,:\,t_{k-1}<\tau_n\le t_k\}} e^m(\tau_n).
\end{equation}
Setting $M=1$ reproduces the original best-level OFI of Cont et al.\ (2014)
exactly. Hidden orders are excluded.

\paragraph{No normalisation.} \textbf{This paper applies none at all} --- not by
depth, event count, spread or volatility. MLOFI is a raw signed share count, and
the paper lists rescaling as future work. Every other construction in this
document divides by something, which matters when comparing coefficient
magnitudes across them.

\subsection{Model and estimation}

\begin{equation}
\Delta P(t_{k-1},t_k)=\alpha+\sum_{m=1}^{M}\beta_m\,\mathrm{MLOFI}^m(t_{k-1},t_k)+\varepsilon,
\qquad M=10,
\end{equation}
with $\Delta P(t,t+\Delta t)=P(t+\Delta t)-P(t)$ the contemporaneous mid-price
change. Ridge minimises $\lVert y-X\beta\rVert_2^2+\lambda\lVert\beta\rVert_2^2$
with $\lambda$ chosen by 5-fold cross-validation over 50 logarithmically spaced
candidates in $[10^{-5},10^{5}]$. The fitted penalties differ by two orders of
magnitude across stocks: AMZN 139.0, TSLA 54.3, NFLX 33.9, ORCL 3.2, CSCO 1.3,
MU 2.0. In the paper's ridge formulation the design matrix carries a leading
column of ones, so the intercept sits inside the penalised vector; no
standardisation of regressors is described.

\paragraph{Data.} LOBSTER, Nasdaq, event-by-event, 4 January to 30 December
2016. Six stocks split by relative tick size: AMZN, TSLA, NFLX (small-tick) and
ORCL, CSCO, MU (large-tick). Sessions 10:00--15:30. Each day is split into
$I=11$ non-overlapping 30-minute windows, each into $K=180$ buckets of 10
seconds, giving one regression per window and $252\times 11=2{,}772$ fits per
stock.

\subsection{Results}

\begin{table}[h]
\centering
\captionsetup{font=small}
\caption{Out-of-sample RMSE in ticks at $M=10$, against best-level OFI.
Improvements are as the paper reports them.}
\small
\begin{tabular}{lcccccc}
\toprule
& AMZN & TSLA & NFLX & ORCL & CSCO & MU \\
\midrule
OFI (best level) & 9.72 & 5.35 & 2.03 & 0.25 & 0.19 & 0.22 \\
MLOFI, OLS       & 8.53 & 4.97 & 1.49 & 0.09 & 0.06 & 0.09 \\
MLOFI, ridge     & 8.05 & 4.53 & 1.41 & 0.08 & 0.05 & 0.08 \\
\midrule
Ridge improvement & 17\% & 15\% & 31\% & 68\% & 74\% & 64\% \\
\bottomrule
\end{tabular}
\end{table}

Depth matters most where the tick is large: roughly 65--75\% RMSE improvement
for large-tick stocks against 15--30\% for small-tick ones. Under OLS the
out-of-sample error for AMZN and TSLA turns \emph{upward} beyond about $M=5$;
under ridge it does not. That single contrast is the argument for regularising.

\begin{lstlisting}[caption={MLOFI with ridge, cross-validated as in the paper.}]
def mlofi(book, levels, bucket):
    """Raw MLOFI per level. No normalisation, per Xu et al."""
    grouper = pd.Grouper(freq=bucket, label="right", closed="right")
    cols = {}
    for m in range(1, levels + 1):
        bp, bs = book[f"bid_px_{m}"], book[f"bid_sz_{m}"]
        ap, asz = book[f"ask_px_{m}"], book[f"ask_sz_{m}"]
        bp0, bs0, ap0, as0 = bp.shift(1), bs.shift(1), ap.shift(1), asz.shift(1)
        dW = np.where(bp > bp0, bs, np.where(bp == bp0, bs - bs0, -bs0))
        dV = np.where(ap > ap0, -as0, np.where(ap == ap0, asz - as0, asz))
        cols[f"mlofi_{m}"] = pd.Series(dW - dV, index=book.index).groupby(grouper).sum()
    return pd.DataFrame(cols)


RIDGE_GRID = np.logspace(-5, 5, 50)          # paper: 50 points in [1e-5, 1e5]

def fit_mlofi_ridge(feat, mid_change, window=180):
    return rolling_fit_predict(
        feat, mid_change, window=window,
        make_model=lambda: RidgeCV(alphas=RIDGE_GRID, cv=5),
    )
\end{lstlisting}

%======================================================================
\section{Integrated OFI and cross-impact}\label{sec:iofi}

\textbf{Source.} Cont, Cucuringu \& Zhang (2023), \emph{Cross-Impact of Order
Flow Imbalance in Equity Markets}, Quantitative Finance 23(10).

\subsection{Definition}

Per-event flows for stock $i$ at level $m$:
\begin{equation}
\mathrm{OF}^{m,b}_{i,n}=
\begin{cases}
q^{m,b}_{i,n}, & P^{m,b}_{i,n}>P^{m,b}_{i,n-1},\\
q^{m,b}_{i,n}-q^{m,b}_{i,n-1}, & P^{m,b}_{i,n}=P^{m,b}_{i,n-1},\\
-q^{m,b}_{i,n}, & P^{m,b}_{i,n}<P^{m,b}_{i,n-1},
\end{cases}
\qquad
\mathrm{OF}^{m,a}_{i,n}=
\begin{cases}
-q^{m,a}_{i,n}, & P^{m,a}_{i,n}>P^{m,a}_{i,n-1},\\
q^{m,a}_{i,n}-q^{m,a}_{i,n-1}, & P^{m,a}_{i,n}=P^{m,a}_{i,n-1},\\
q^{m,a}_{i,n}, & P^{m,a}_{i,n}<P^{m,a}_{i,n-1}.
\end{cases}
\end{equation}
Aggregating and scaling by average depth:
\begin{equation}
\mathrm{OFI}^{m,h}_{i,t}=\sum_{n=N(t-h)+1}^{N(t)}\bigl(\mathrm{OF}^{m,b}_{i,n}-\mathrm{OF}^{m,a}_{i,n}\bigr),
\qquad
\mathrm{ofi}^{m,h}_{i,t}=\frac{\mathrm{OFI}^{m,h}_{i,t}}{Q^{M,h}_{i,t}},
\end{equation}
\begin{equation}
Q^{M,h}_{i,t}=\frac{1}{M}\sum_{m=1}^{M}\frac{1}{2\,\Delta N(t)}
\sum_{n=N(t-h)+1}^{N(t)}\bigl(q^{m,b}_{i,n}+q^{m,a}_{i,n}\bigr).
\end{equation}
One \emph{common} denominator scales every level, which preserves the relative
size of the levels. The motivation is the intraday pattern in depth: without it
the same raw flow means different things at 10:00 and at 15:00.

The ten scaled levels are then compressed into one number:
\begin{equation}
\mathrm{ofi}^{I,h}_{i,t}=\frac{\mathbf{w}_1^{\!\top}\,\mathbf{ofi}^{(h)}_{i,t}}{\lVert\mathbf{w}_1\rVert_1},
\qquad
\mathbf{ofi}^{(h)}_{i,t}=\bigl(\mathrm{ofi}^{1,h}_{i,t},\dots,\mathrm{ofi}^{10,h}_{i,t}\bigr)^{\!\top}.
\end{equation}
The first principal component explains 89.06\% (sd 6.12) of the variance of the
level vector; a plain average across levels reaches 85.07\%.

\paragraph{Under-specified step.} The paper says only that $\mathbf{w}_1$ is
``the first principal vector computed from historical data''. It does not state
whether the PCA is on the correlation or the covariance matrix, nor the
estimation window. Results are reported ``averaged across each stock and
trading day'', which suggests a per-stock, per-day fit, but the paper never
states that as the rule. Fix your own choice and record it.

\subsection{Models}

With log mid-price returns $r^{(h)}_{i,t}=\log\bigl(P_{i,t}/P_{i,t-h}\bigr)$:
\begin{align}
\text{PI}^{[1]}&:\ r^{(h)}_{i,t}=\alpha_i+\beta_i\,\mathrm{ofi}^{1,h}_{i,t}+\epsilon_{i,t},
&\text{PI}^{I}&:\ r^{(h)}_{i,t}=\alpha_i+\beta_i\,\mathrm{ofi}^{I,h}_{i,t}+\epsilon_{i,t},\\
\text{CI}^{[1]}&:\ r^{(h)}_{i,t}=\alpha_i+\beta_{i,i}\,\mathrm{ofi}^{1,h}_{i,t}
+\sum_{j\neq i}\beta_{i,j}\,\mathrm{ofi}^{1,h}_{j,t}+\eta_{i,t},
&\text{CI}^{I}&:\ \text{same with } \mathrm{ofi}^{I,h}.
\end{align}
The forecasting versions predict the next $f$-minute return from lagged
aggregates over $kh$, $k\in L=\{1,2,3,5,10,20,30\}$:
\begin{align}
\text{FPI}^{[1]}&:\ r^{(f)}_{i,t+f}=\alpha_i+\sum_{k\in L}\beta^{k}_i\,\mathrm{ofi}^{1,(kh)}_{i,t}+\epsilon_{i,t+f},\\
\text{FCI}^{[1]}&:\ r^{(f)}_{i,t+f}=\alpha_i+\sum_{j=1}^{N}\sum_{k\in L}\beta^{k}_{i,j}\,\mathrm{ofi}^{1,(kh)}_{j,t}+\eta_{i,t+f},
\end{align}
with $\mathrm{ofi}^{I,(kh)}$ in place of $\mathrm{ofi}^{1,(kh)}$ for the
integrated variants FPI$^{I}$ and FCI$^{I}$. The benchmarks are AR (own lagged
returns) and CAR (all stocks' lagged returns), with the same lag set.

\paragraph{Estimation.} PI and AR by OLS; CI, FCI and CAR by LASSO with a
cross-validated $\ell_1$ penalty (no fold count or grid stated, and no
standardisation stated). LASSO is required because $N\approx100$ regressors
face 30 one-minute observations per window. Non-overlapping 30-minute windows
within 10:00--15:30; one-minute buckets; physical time, not trade time; Nasdaq
ITCH via LOBSTER, top 100 S\&P 500 names, 2017--2019.

\subsection{Results}

\begin{table}[h]
\centering
\captionsetup{font=small}
\caption{Contemporaneous $R^2$ in percentage points, and one-minute-ahead
out-of-sample $R^2$. Cross-impact adds 2.71 points in sample on best-level OFI
but only 0.71 on integrated OFI, and nothing out of sample.}
\small
\begin{tabular}{lcccc}
\toprule
& PI$^{[1]}$ & CI$^{[1]}$ & PI$^{I}$ & CI$^{I}$ \\
\midrule
In-sample, contemporaneous & 71.16 & 73.87 & 87.14 & 87.85 \\
Out-of-sample, contemporaneous & 64.64 & 66.03 & 83.83 & 83.62 \\
Out-of-sample, one minute ahead & $-0.37$ & $-0.10$ & $-0.36$ & $-0.10$ \\
\bottomrule
\end{tabular}
\end{table}

For forecasting, the cross-sectional models beat their own-stock counterparts
(AR $-0.36$, CAR $-0.10$), but integration brings no advantage: the PCA step
discards which level the flow came from, and level identity carries predictive
content. Annualised PnL under the paper's trading rule, ignoring costs: 0.21
(FPI$^{[1]}$), 0.43 (FCI$^{[1]}$), 0.23 (FPI$^{I}$), 0.39 (FCI$^{I}$), 0.23
(AR), 0.40 (CAR).

The trading rule itself is worth stating, because both this paper and Sitaru et
al.\ evaluate with it:
\begin{equation}
w_{i,t}=\frac{\ind{|f_{i,t}|>\mathrm{sprd}_{i,t}}\cdot f_{i,t}/\sigma_{i,t}}
{\sum_{n=1}^{N}\ind{|f_{n,t}|>\mathrm{sprd}_{n,t}}\cdot f_{n,t}/\sigma_{n,t}},
\end{equation}
where $f_{i,t}$ is the model's forecast for the next minute,
$\mathrm{sprd}_{i,t}$ the relative bid-ask spread, and $\sigma_{i,t}$ the
standard deviation of the model's in-sample forecasts over the previous 30
minutes. If no forecast exceeds its spread, nothing trades that minute.

\begin{lstlisting}[caption={Integrated OFI. The PCA is fitted on past data only; the
window is a choice the paper leaves open, so it is a parameter here.}]
def depth_scaled_ofi(book, levels, bucket):
    """Scaled level OFIs sharing one depth denominator Q^{M,h}."""
    grouper = pd.Grouper(freq=bucket, label="right", closed="right")
    n_events = book.groupby(grouper).size()
    depth = []
    for m in range(1, levels + 1):
        depth.append((book[f"bid_sz_{m}"] + book[f"ask_sz_{m}"]).groupby(grouper).sum())
    Q = pd.concat(depth, axis=1).sum(axis=1) / (2.0 * levels * n_events)

    cols = {}
    for m in range(1, levels + 1):
        bp, bs = book[f"bid_px_{m}"], book[f"bid_sz_{m}"]
        ap, asz = book[f"ask_px_{m}"], book[f"ask_sz_{m}"]
        bp0, bs0, ap0, as0 = bp.shift(1), bs.shift(1), ap.shift(1), asz.shift(1)
        ofb = np.where(bp > bp0, bs, np.where(bp == bp0, bs - bs0, -bs))
        ofa = np.where(ap > ap0, -asz, np.where(ap == ap0, asz - as0, asz))
        raw = pd.Series(ofb - ofa, index=book.index).groupby(grouper).sum()
        cols[f"ofi_{m}"] = raw / Q
    return pd.DataFrame(cols)


from sklearn.decomposition import PCA

def integrated_ofi(scaled, fit_window):
    """Project onto the first PC of past scaled OFIs, l1-normalised weights.

    fit_window rows ending before t are used, so the weight vector applied at t
    never sees t. The paper does not state the window; record whatever you pick.
    """
    out = pd.Series(index=scaled.index, dtype=float)
    vals = scaled.to_numpy()
    for t in range(fit_window, len(scaled)):
        past = vals[t - fit_window:t]
        w = PCA(n_components=1).fit(past).components_[0]
        out.iloc[t] = float(vals[t] @ w / np.abs(w).sum())
    return out
\end{lstlisting}

%======================================================================
\section{Decomposed OFI by event type}\label{sec:dofi}

\textbf{Source.} Sitaru, Calinescu \& Cucuringu (2023), \emph{Order Flow
Decomposition for Price Impact Analysis in Equity Limit Order Books}, ICAIF '23.

\subsection{Definition}

For event $n$ at level $\ell$,
\begin{equation}
\mathrm{aOF}_n^{\ell}=
\begin{cases}
-\,qa_{n-1}^{\ell}, & a_n^{\ell}>a_{n-1}^{\ell},\\
qa_n^{\ell}-qa_{n-1}^{\ell}, & a_n^{\ell}=a_{n-1}^{\ell},\\
qa_n^{\ell}, & \text{otherwise},
\end{cases}
\qquad
\mathrm{bOF}_n^{\ell}=
\begin{cases}
qb_n^{\ell}, & b_n^{\ell}>b_{n-1}^{\ell},\\
qb_n^{\ell}-qb_{n-1}^{\ell}, & b_n^{\ell}=b_{n-1}^{\ell},\\
-\,qb_{n-1}^{\ell}, & \text{otherwise},
\end{cases}
\end{equation}
\begin{equation}
\mathrm{OFI}_t^{\ell,h}=\sum_{n=N(t-h)+1}^{N(t)}\bigl(\mathrm{bOF}_n^{\ell}-\mathrm{aOF}_n^{\ell}\bigr),
\qquad
Q_t^{L,h}=\frac{1}{L}\sum_{\ell=1}^{L}\frac{1}{\Delta N(t)}
\sum_{n=N(t-h)+1}^{N(t)}\frac{qa_n^{\ell}+qb_n^{\ell}}{2},
\end{equation}
and $\mathrm{ofi}_t^{\ell,h}=\mathrm{OFI}_t^{\ell,h}/Q_t^{L,h}$.

The decomposition splits the same sum by event type:
\begin{equation}
\mathrm{OFIa}_t^{\ell,h}=\sum_{n=N(t-h)+1}^{N(t)}
\bigl(\mathrm{bOF}_n^{\ell}-\mathrm{aOF}_n^{\ell}\bigr)\ind{\mathrm{type}(n)=\mathrm{add}},
\end{equation}
with $\mathrm{OFIc}$ (cancel) and $\mathrm{OFIt}$ (trade) analogous, so that
\begin{equation}
\mathrm{OFI}_t^{\ell,h}=\mathrm{OFIa}_t^{\ell,h}+\mathrm{OFIc}_t^{\ell,h}+\mathrm{OFIt}_t^{\ell,h}.
\end{equation}
Each component is divided by the same $Q_t^{L,h}$, so additivity survives
normalisation. Event types need market-by-order data; they cannot be recovered
from snapshots.

\paragraph{LOBSTER mapping.} add $=$ type 1; cancel $=$ types 2 and 3; trade $=$
types 4, 5 and 6. Execution of a hidden order does not change the book, so its
contribution is null.

\subsection{Models and estimation}

\begin{align}
\mathrm{PI}^{\ell}&:\ r_t^h=\alpha+\sum_{k=1}^{\ell}\beta_k\,\mathrm{ofi}_t^{k,h}+\epsilon,\\
\mathrm{PID}^{\ell}&:\ r_t^h=\alpha+\sum_{k=1}^{\ell}
\bigl(\beta_{k,\mathrm{a}}\mathrm{ofia}_t^{k,h}+\beta_{k,\mathrm{c}}\mathrm{ofic}_t^{k,h}
+\beta_{k,\mathrm{t}}\mathrm{ofit}_t^{k,h}\bigr)+\epsilon,
\end{align}
and forecasting versions $\mathrm{FPI}^{\ell}_M$, $\mathrm{FPID}^{\ell}_M$ with
lag sets $M_3=\{1,2,3\}$ and $M_7=\{1,2,3,5,10,20,30\}$ minutes. Contemporaneous:
$h=10$ s, 10 levels, 180-observation window, next 30 minutes out of sample.
Forecasting: $h=f=1$ min, rolling 30-observation window, LASSO with
cross-validation except the two best-level OLS baselines. Same LOBSTER universe
as Cont et al.

\subsection{Results}

\begin{table}[h]
\centering
\captionsetup{font=small}
\caption{Contemporaneous adjusted $R^2$, percentage points (sd in parentheses).
Out-of-sample values for decomposed OFI under OLS at $\ell\ge6$ fall below
$-100\%$ and are not reported by the authors.}
\small
\begin{tabular}{lcc}
\toprule
Model & In-sample & Out-of-sample \\
\midrule
$\mathrm{PI}^{1}$ (OLS)  & 73.67 (14.35) & 69.42 (22.31) \\
$\mathrm{PI}^{10}$ (OLS) & 91.97 (8.73)  & 89.37 (21.06) \\
$\mathrm{PID}^{10}$ (OLS) & 93.00 (7.37) & not reported \\
$\mathrm{PID}^{10}$ (LASSO) & 92.23 (8.82) & 87.03 (14.85) \\
$\mathrm{PID}^{I,10}$ (10 components) & 90.57 (11.01) & 86.24 (19.34) \\
\bottomrule
\end{tabular}
\end{table}

Contemporaneously the decomposition adds nothing: about one point in sample,
about two given back out of sample. The principal-component structure differs
sharply from section~\ref{sec:iofi}: the first component of the 30-variable
decomposed vector explains only 59.43\% of its variance, and five components are
needed for 89.24\%. Event type genuinely adds dimensions.

Forecasting is where it pays. Every out-of-sample $R^2$ is negative --- the best
is $\mathrm{FPI}^1_{M_3}$ at $-0.082$ --- while the trading evaluation improves
several-fold:

\begin{table}[h]
\centering
\captionsetup{font=small}
\caption{Forecast-implied PnL in basis points per dollar traded, and Sharpe
ratio, under the trading rule of section~\ref{sec:iofi}. Costs excluded.}
\small
\begin{tabular}{lccccc}
\toprule
& FPI$^1$ (OLS) & FPI$^1$ (LASSO) & FPID$^1$ & FPI$^{10}$ & FPID$^{10}$ \\
\midrule
$M_3$ & 0.22 / 1.24 & 0.59 / 2.17 & 1.01 / 4.26 & 0.54 / 2.09 & \textbf{1.66 / 6.90} \\
$M_7$ & 0.68 / 5.86 & 1.26 / 5.22 & 2.42 / 10.65 & 1.33 / 5.18 & \textbf{3.15 / 13.51} \\
\bottomrule
\end{tabular}
\end{table}

Writing $\mathcal{F}(\cdot)$ for the LASSO selection frequency, at every level
and lag
\[
\mathcal{F}(\mathrm{adds})>\mathcal{F}(\mathrm{cancels})>\mathcal{F}(\mathrm{trades}),
\]
and in the forecasting model the trade component is rarely selected at all.
Order submissions, not executions, carry the forecasting content. Caveat from
the authors: $\mathrm{FPID}^{10}_{M_7}$ has 210 regressors against 30
observations, and 53\% of its regressions select nothing but the intercept, so
the headline Sharpe ratio rests on a minority of active minutes.

\begin{lstlisting}[caption={Decomposed OFI. The assertion enforces the additivity identity,
which is the cheapest test that the event mapping is right.}]
ADD, CANCEL, TRADE = "add", "cancel", "trade"
LOBSTER_TYPE_MAP = {1: ADD, 2: CANCEL, 3: CANCEL, 4: TRADE, 5: TRADE, 6: TRADE}


def decomposed_ofi(book, levels, bucket):
    """Normalised OFI and its add/cancel/trade parts, per level and bucket."""
    etype = book["event_type"]                       # already mapped
    grouper = pd.Grouper(freq=bucket, label="right", closed="right")
    n_events = book.groupby(grouper).size()

    per_level = [((book[f"ask_sz_{l}"] + book[f"bid_sz_{l}"]) / 2.0)
                 .groupby(grouper).sum() for l in range(1, levels + 1)]
    Q = pd.concat(per_level, axis=1).sum(axis=1) / (levels * n_events)

    cols = {}
    for l in range(1, levels + 1):
        bp, bs = book[f"bid_px_{l}"], book[f"bid_sz_{l}"]
        ap, asz = book[f"ask_px_{l}"], book[f"ask_sz_{l}"]
        bp0, bs0, ap0, as0 = bp.shift(1), bs.shift(1), ap.shift(1), asz.shift(1)
        bof = np.where(bp > bp0, bs, np.where(bp == bp0, bs - bs0, -bs0))
        aof = np.where(ap > ap0, -as0, np.where(ap == ap0, asz - as0, asz))
        signed = pd.Series(bof - aof, index=book.index).fillna(0.0)
        cols[f"ofi_{l}"] = signed.groupby(grouper).sum() / Q
        for name in (ADD, CANCEL, TRADE):
            cols[f"ofi_{name}_{l}"] = signed.where(etype == name, 0.0).groupby(grouper).sum() / Q

    out = pd.DataFrame(cols).dropna(how="all")
    for l in range(1, levels + 1):
        parts = out[[f"ofi_{n}_{l}" for n in (ADD, CANCEL, TRADE)]].sum(axis=1)
        assert np.allclose(parts, out[f"ofi_{l}"], atol=1e-8, equal_nan=True), l
    return out


def lagged_blocks(feat, lags):
    """Aggregate each feature over m-bucket windows, one block per lag."""
    frames = []
    for m in lags:
        rolled = feat.rolling(window=m, min_periods=m).sum()
        rolled.columns = [f"{c}_lag{m}" for c in feat.columns]
        frames.append(rolled)
    return pd.concat(frames, axis=1)


M3, M7 = (1, 2, 3), (1, 2, 3, 5, 10, 20, 30)

def fit_forecast(feat, fwd_ret, lags=M7, window=30):
    X = lagged_blocks(feat, lags).dropna()
    y = fwd_ret.reindex(X.index)
    return rolling_fit_predict(
        X, y, window=window,
        make_model=lambda: make_pipeline(StandardScaler(),
                                         LassoCV(cv=3, n_alphas=25, max_iter=5000)),
    )
\end{lstlisting}

%======================================================================
\section{Generalised OFI for snapshot data}\label{sec:gofi}

\textbf{Source.} Su, Sun, Li \& Yuan (2021), \emph{The Price Impact of
Generalized Order Flow Imbalance}, arXiv:2112.02947.

\subsection{The problem and the fix}

Standard OFI accounts for changes at the \emph{position} of the best quote,
which is exact only if the quote moves at most one tick between observations.
Chinese exchanges publish snapshots every three seconds, over which the touch
can traverse several levels. Generalised OFI keys on the \emph{value} of the
best price and sums the quantities at every level traversed:
\begin{equation}
\mathrm{GOFI}_{k-1,k}=\sum_{n=1}^{N} gene_n, \qquad gene_n=\Delta GBid_n-\Delta GAsk_n,
\end{equation}
\begin{equation}
\Delta GBid_n=
\begin{cases}
\sum_{i\le m^b} v^b_{n,i}-v^b_{n-1,1}, & P^b_n>P^b_{n-1},\\
v^b_{n,1}-v^b_{n-1,1}, & P^b_n=P^b_{n-1},\\
v^b_{n,1}-\sum_{i\le m^b} v^b_{n-1,i}, & P^b_n<P^b_{n-1},
\end{cases}
\qquad
\Delta GAsk_n=
\begin{cases}
v^a_{n,1}-\sum_{i\le m^a} v^a_{n-1,i}, & P^a_n>P^a_{n-1},\\
v^a_{n,1}-v^a_{n-1,1}, & P^a_n=P^a_{n-1},\\
\sum_{i\le m^a} v^a_{n,i}-v^a_{n-1,1}, & P^a_n<P^a_{n-1},
\end{cases}
\end{equation}
where $v^{b}_{n,i}$ is the queue size at the $i$-th bid level at $n$, and the
number of levels crossed is
\begin{equation}
m^b=\left|\frac{P^b_n-P^b_{n-1}}{\delta}\right|+1,
\qquad
m^a=\left|\frac{P^a_n-P^a_{n-1}}{\delta}\right|+1,
\end{equation}
with $\delta$ the tick size.

\paragraph{Stationarised variants.} log-OFI and log-GOFI apply the logarithm to
the \emph{queue quantities} before differencing, not to the imbalance. The paper
states this in words only and prints no display equation for log-OFI, so do not
present one as theirs. The motivation is that queue sizes vary enormously
across levels, which the fixed-depth assumption of the original model does not
accommodate.

\subsection{Results}

Univariate regressions of the mid-price change on each indicator, at horizons of
30 seconds, 1 minute and 5 minutes; ten CSI 500 constituents; three-second
snapshots; in-sample January to March 2021, out-of-sample April to June 2021.
The paper gives no formula for how its $R^2$ values are computed, so fix and
record your own (the harness in section~\ref{sec:protocol} uses a zero
benchmark).

\begin{table}[h]
\centering
\captionsetup{font=small}
\caption{Average $R^2$ across the ten stocks, percentage points, in-sample /
out-of-sample.}
\small
\begin{tabular}{lcccc}
\toprule
Horizon & OFI & log-OFI & GOFI & log-GOFI \\
\midrule
30 seconds & 38.38 / 32.89 & 75.96 / 76.37 & 45.73 / 40.35 & \textbf{83.58 / 83.57} \\
1 minute   & 44.45 / 38.13 & 77.63 / 77.85 & 51.88 / 46.27 & \textbf{85.46 / 85.37} \\
5 minutes  & 49.64 / 42.57 & 77.04 / 77.36 & 56.76 / 51.65 & \textbf{85.76 / 86.01} \\
\bottomrule
\end{tabular}
\end{table}

Three things stand out. The log transform buys more than the generalisation
does, contrary to the paper's framing. The two are close to additive. And for
the generalised indicators the in-sample and out-of-sample values essentially
coincide, where plain OFI loses five to seven points.

\paragraph{Read this one carefully.} GOFI aggregates \emph{resting depth across
traversed levels} rather than incremental flow at a fixed level, so it is closer
to a state variable of the book than to a flow variable. The paper does not
discuss whether that mechanically narrows the gap to the contemporaneous
mid-price change. Anyone building on it should check that separately.

\begin{lstlisting}[caption={Generalised OFI from snapshots, with the optional log transform.}]
def generalized_ofi(snaps, tick, n_levels=5, use_log=False):
    """GOFI / log-GOFI from L2 snapshots (Su et al.).

    `snaps` has bid_px_1..n, bid_sz_1..n, ask_px_1..n, ask_sz_1..n per snapshot.
    """
    def v(side, i, shift=0):
        s = snaps[f"{side}_sz_{i}"]
        s = s.shift(shift) if shift else s
        return np.log1p(s) if use_log else s

    bp, ap = snaps["bid_px_1"], snaps["ask_px_1"]
    bp0, ap0 = bp.shift(1), ap.shift(1)
    mb = (bp - bp0).abs().div(tick).round().astype("Int64").fillna(0) + 1
    ma = (ap - ap0).abs().div(tick).round().astype("Int64").fillna(0) + 1

    def cum(side, upto, shift):
        """Sum of the first `upto` levels, row-wise, capped at n_levels."""
        acc = pd.Series(0.0, index=snaps.index)
        for i in range(1, n_levels + 1):
            acc = acc.add(v(side, i, shift).where(upto >= i, 0.0), fill_value=0.0)
        return acc

    d_bid = np.where(bp > bp0, cum("bid", mb, 0) - v("bid", 1, 1),
             np.where(bp == bp0, v("bid", 1, 0) - v("bid", 1, 1),
                      v("bid", 1, 0) - cum("bid", mb, 1)))
    d_ask = np.where(ap > ap0, v("ask", 1, 0) - cum("ask", ma, 1),
             np.where(ap == ap0, v("ask", 1, 0) - v("ask", 1, 1),
                      cum("ask", ma, 0) - v("ask", 1, 1)))
    return pd.Series(d_bid - d_ask, index=snaps.index, name="gofi")
\end{lstlisting}

%======================================================================
\section{OFI as a shock: horizon choice and screening}\label{sec:hu}

\textbf{Source.} Hu \& Zhang (2025), \emph{Stochastic Price Dynamics in Response
to Order Flow Imbalance: Evidence from CSI 300 Index Futures}, arXiv:2505.17388.

\subsection{Four metrics, stated in one place}

\begin{equation}
\mathrm{OFI}_k=\sum_{n=N(t_{k-1})+1}^{N(t_k)} e_n,
\qquad
\mathrm{TI}_k=\sum_{n=N(t_{k-1})+1}^{N(t_k)} \omega_n,
\end{equation}
\begin{align}
e_n&=\ind{p^B_n\ge p^B_{n-1}}q^B_n-\ind{p^B_n\le p^B_{n-1}}q^B_{n-1}
-\ind{p^A_n\le p^A_{n-1}}q^A_{n-1}+\ind{p^A_n\ge p^A_{n-1}}q^A_n,\\
\omega_n&=\Bigl(\ind{p^{last}_n>M_p}-\ind{p^{last}_n<M_p}\notag\\
&\qquad\quad+\ind{p^{last}_n=M_p,\ p^{last}_n\ge p^{last}_{n-1}}
-\ind{p^{last}_n=M_p,\ p^{last}_n<p^{last}_{n-1}}\Bigr)v,
\end{align}
with $M_p=(P^A_n+P^B_n)/2$ and $v=V_n-V_{n-1}$ the incremental volume. The
indicator form of $e_n$ is the same object as the piecewise definitions above,
written compactly.
\begin{equation}
\mathrm{Lambda}_k=\frac{p^{max}_k-p^{min}_k}{N(t_k)-N(t_{k-1})},
\qquad
\mathrm{AvgEn}_k=\mathrm{OFI}^{avg}_{N(t_k)}-\mathrm{OFI}^{avg}_{N(t_{k-1})},
\quad
\mathrm{OFI}^{avg}_{N(t_k)}=\frac{1}{N(t_k)}\sum_{n=1}^{N(t_k)}e_n.
\end{equation}

\subsection{The shock model}

The point of the paper is that a fixed window and a matching forecast horizon
are an arbitrary pairing. An OFI accumulated over a window is instead treated as
a step shock whose response is measured across many horizons. The drift of
geometric Brownian motion is replaced by a mean-reverting process driven by a
jump L\'evy process:
\begin{equation}
dS_t=\mu_t S_t\,dt+\sigma S_t\,dW_t,
\qquad
d\mu_t=\theta(\mu_l-\mu_t)\,dt+dL_t,
\qquad \mu_0=\mathrm{OFI}_0\,\rho\,k,
\end{equation}
with $\mu_l$ the long-run drift (set to zero empirically), $\theta$ the
mean-reversion speed and $L_t$ a zero-mean symmetric L\'evy process with
$\sigma_L^2=\E[L_t^2]/t$. The L\'evy driver is an empirical choice: tick-level
flow contributions are heavy-tailed, so a Gaussian driver misfits. For
$R_t=\log(S_t/S_0)$,
\begin{equation}
\E[R_t]=\mu_0\frac{1-e^{-\theta t}}{\theta}-\frac{\sigma^2}{2}t,
\qquad
\Var(R_t)=\sigma^2 t+\frac{\sigma_L^2}{\theta^2}
\Bigl[t-\frac{2}{\theta}\bigl(1-e^{-\theta t}\bigr)+\frac{1}{2\theta}\bigl(1-e^{-2\theta t}\bigr)\Bigr],
\end{equation}
and the \emph{response ratio} (their quasi-Sharpe ratio) is the ratio of the two,
\begin{equation}
QS(t)=\frac{\E[R_t]}{\sqrt{\Var(R_t)}} .
\end{equation}
It vanishes as $t\to0^+$ like $\sqrt{t}$ and tends to a constant as
$t\to\infty$. Practically it says where in horizon space a signal is worth
trading: the numerator saturates while the denominator keeps growing.

\subsection{A screening test worth reusing}

The paper's most portable contribution is an ex-ante test for keeping any
microstructure indicator. A robust metric should
\begin{enumerate}[leftmargin=2em,itemsep=1pt]
\item correlate significantly with contemporaneous price moves, with
coefficients stable across time windows;
\item show memory --- high autocorrelation with slow decay;
\item preserve both properties across market regimes, varying quantitatively
but \emph{never in sign}.
\end{enumerate}
Sign reversal across regimes is the disqualifying failure.

\paragraph{Data and grid.} CSI 300 index futures; exchange snapshots every
500\,ms, so one ``tick'' is 500\,ms; roughly six million ticks over
December 2023 to November 2024. Historical windows of 1, 2, 5, 10, 20, 30, 60,
120 and 240 ticks are crossed against forecast horizons of 1 to 3600 ticks
(30 minutes). LASSO with 5-fold cross-validation, in-sample to out-of-sample
ratio 4:1.

\begin{lstlisting}[caption={Response-ratio curve and the regime-stability screen.}]
def response_ratio(t, mu0, theta, sigma, sigma_L):
    """QS(t) for a step OFI shock (Hu & Zhang, eq. 24)."""
    t = np.asarray(t, dtype=float)
    mean = mu0 * (1.0 - np.exp(-theta * t)) / theta - 0.5 * sigma**2 * t
    var = sigma**2 * t + (sigma_L**2 / theta**2) * (
        t - (2.0 / theta) * (1.0 - np.exp(-theta * t))
        + (1.0 / (2.0 * theta)) * (1.0 - np.exp(-2.0 * theta * t)))
    return mean / np.sqrt(var)


def screen_indicator(feature, ret, regime_labels, max_lag=120):
    """The three criteria, as numbers you can threshold.

    Returns contemporaneous correlation, its dispersion across regimes, the
    autocorrelation decay, and whether the sign is stable across regimes.
    """
    corr_all = feature.corr(ret)
    by_regime = {r: feature[regime_labels == r].corr(ret[regime_labels == r])
                 for r in regime_labels.unique()}
    acf = [feature.autocorr(lag) for lag in range(1, max_lag + 1)]
    signs = {np.sign(v) for v in by_regime.values() if np.isfinite(v)}
    return {
        "corr": corr_all,
        "corr_by_regime": by_regime,
        "corr_dispersion": float(np.nanstd(list(by_regime.values()))),
        "acf_halflife": next((i for i, a in enumerate(acf, 1) if a < acf[0] / 2), np.nan),
        "sign_stable": len(signs) == 1,          # criterion 3: the hard one
    }
\end{lstlisting}

%======================================================================
\section{Generalised imbalance from trades alone}\label{sec:genimb}

\textbf{Sources.} Maitrier \& Bouchaud (2025), arXiv:2506.07711, and the
simulation companion Maitrier, Loeper \& Bouchaud (2025), arXiv:2509.05065.

This is the construction that ports to markets without a visible book, because
it needs only signed trades and their sizes.

\subsection{Definition}

\begin{equation}
I^a_T=\sum_{t\in[0,T]}\varepsilon_t\,(q_t)^a .
\end{equation}
The exponent $a$ sets how much large trades count: $a=0$ is pure sign
imbalance, $a=1$ is volume imbalance, and intermediate values interpolate. The
companion paper writes the same object in integral form,
$I^a_T=\int_0^T dt\,\varepsilon_t q_t^a$.

The framework around it starts from the square-root law,
\begin{equation}
I(Q)=Y\,\sigma\sqrt{\frac{Q}{\varphi}},
\end{equation}
with $Q$ the metaorder volume, $\sigma$ the volatility, $Y$ an $O(1)$ constant
and $\varphi$ the average executed order flow per unit time. Metaorders arrive
at rate $\nu$ with signs $\varepsilon$, durations drawn from
$\Psi_q(s)=\mu_q s_0^{\mu_q}/s^{1+\mu_q}$ and child orders executed at rate
$\phi$. Two dependencies do the work:
\begin{equation}
\mu_q=\mu_1+\lambda\ell,
\qquad
\beta_q=\beta_1-\lambda'\ell,
\qquad \ell=\log q,
\end{equation}
so larger child orders are less persistent and their impact decays more slowly,
and a long-range correlation \emph{between} metaorder signs,
$\E[\varepsilon(t)\varepsilon(t+\tau)]\simeq\Gamma(\tau_0/\tau)^{\gamma_\times}$.

\paragraph{Diffusion condition.} Transient impact with heavy-tailed durations
alone gives sub-diffusion. Diffusive prices require
\begin{equation}
\gamma_\times=1-2\beta,
\end{equation}
which with $\beta=0.2$ gives $\gamma_\times=0.6$. This is the paper's central
falsifiable claim.

\subsection{What to measure on data}

The testable object is the correlation between returns and generalised
imbalance,
\begin{equation}
R_a(T)=\frac{\E[\Delta_T\,I^a_T]}{\sqrt{\E[\Delta_T^2]}\,\sqrt{\E[(I^a_T)^2]}},
\end{equation}
as a function of $a$ and $T$. Two predictions follow: the scaling exponent of
$\E[\Delta_T I^a_T]$ against $T$ is \emph{non-monotonic} in $a$, and $R_a(T)$ is
\emph{hump-shaped} in $a$ --- there is an optimal weighting of large trades.
Both are observed. Empirically the peak sits at $a^\star\approx0.5$ for LLOY,
TSCO and SPMINI and $a^\star\approx1$ for EUROSTOXX, with peak correlations
around 0.45 for the stocks, 0.6 for SPMINI and 0.75 for EUROSTOXX.

If a large share of volatility came from fundamental news instead, the
covariance would be $T$-independent and free of $a$-dependence. It is not. That
is the paper's argument against the efficient-market reading.

\paragraph{Practical notes.} Everything is measured in \emph{trade time}: $T$
counts trades, not seconds. Translating to real time is non-trivial because
activity is intermittent and has a U-shaped intraday profile; the paper
neglects both. Four assets: LLOY and TSCO (LSE, 2012--2015), SPMINI (2022) and
EUROSTOXX (2016--2018). Trades beyond 1\% of daily volume are clipped. On EUROSTOXX the
sign imbalance collapses onto a master curve with $\chi=0.72$.

\subsection{The simulator, and why it matters for bonds}

The companion paper simulates the model without the analytical approximations
and recovers the same phenomenology. Its most useful result for anyone without
metaorder labels: \textbf{proxy metaorders reconstructed from anonymous trades
reproduce the square-root law, including its prefactor $Y$}. Two conditions are
essential --- preserving the exact trade history and sampling trades without
replacement, and an inter-time threshold so that a long gap between child
orders starts a new metaorder. Agreement degrades toward linear at small
$Q/V_D$, because short proxy metaorders miss the true start, where impact is
most concave. Code: \url{https://github.com/glatouille/ArtificialMarketSimulator.git}.

\begin{lstlisting}[caption={Generalised imbalance and the hump in $a$. This runs on a
trade tape alone: signs, sizes, prices.}]
def generalized_imbalance(signs, sizes, a, window):
    """I^a_T over a rolling window of `window` trades (trade time)."""
    contrib = np.sign(signs) * np.power(np.asarray(sizes, dtype=float), a)
    return pd.Series(contrib, index=signs.index).rolling(window).sum()


def r_a_curve(signs, sizes, log_mid, T, a_grid=np.linspace(0.0, 2.0, 21)):
    """R_a(T) against a. The peak is the optimal weighting of large trades."""
    fwd = log_mid.shift(-T) - log_mid                 # return over T trades
    rows = []
    for a in a_grid:
        imb = generalized_imbalance(signs, sizes, a, T)
        both = pd.concat([fwd, imb], axis=1).dropna()
        num = (both.iloc[:, 0] * both.iloc[:, 1]).mean()
        den = np.sqrt((both.iloc[:, 0] ** 2).mean() * (both.iloc[:, 1] ** 2).mean())
        rows.append({"a": a, "R_a": num / den if den > 0 else np.nan})
    out = pd.DataFrame(rows)
    return out, float(out.loc[out["R_a"].idxmax(), "a"])   # a*, the hump location


def proxy_metaorders(trades, phi, C=4.0):
    """Group anonymous trades into proxy metaorders, per the companion paper.

    Conditions: same sign, no reuse of a trade, and a gap threshold C/phi
    between consecutive child orders.
    """
    out, current, last_t = [], [], None
    gap = C / phi
    for t, row in trades.iterrows():
        if current and (row["sign"] != current[-1]["sign"] or (t - last_t) > gap):
            out.append(current)
            current = []
        current.append({"t": t, **row.to_dict()})
        last_t = t
    if current:
        out.append(current)
    return out
\end{lstlisting}

%======================================================================
\section{Building flow features by clustering}\label{sec:petit}

\textbf{Source.} Petit, Cucuringu \& Cartea (2025), \emph{Data-Driven Trade Flow
Decomposition for Exchange-Traded Funds and their Constituents}, ICAIF '25.

This is the template for turning raw flow into features that a downstream model
can use. The pipeline is: build features around each trade, normalise them
against comparable history, reduce dimension, cluster, then form
cluster-conditional imbalances.

\subsection{The 16 features}

Around a central trade $(i,n)$, count the trades and sum the notional in a
window $\pm\tau$, split three ways: same or opposite direction, ETF or
constituents, before or after. That gives eight cells, each yielding a count and
a notional.

\begin{table}[h]
\centering
\captionsetup{font=small}
\caption{Feature indices. Each cell contributes a trade count and a notional.}
\small
\begin{tabular}{llcc}
\toprule
Direction & Market & Count & Notional \\
\midrule
Same & ETF & 1 & 2 \\
Same & ETF (after) & 3 & 4 \\
Same & Constituents & 5 & 6 \\
Same & Constituents (after) & 7 & 8 \\
Opposite & ETF & 9 & 10 \\
Opposite & ETF (after) & 11 & 12 \\
Opposite & Constituents & 13 & 14 \\
Opposite & Constituents (after) & 15 & 16 \\
\bottomrule
\end{tabular}
\end{table}

All sixteen share one shape. For a market set $\mathcal{S}$ (the ETF
$\mathcal{E}$ or its constituents $\mathcal{M}$), a direction condition $D$
(same or opposite to the central trade) and a window side $W$ (before or after),
\begin{equation}
x_{i,n}=\sum_{j,m} w_{j,m}\cdot
[\![\,j\in\mathcal{S}\,]\!]\cdot
[\![\,D(\varepsilon_{j,m},\varepsilon_{i,n})\,]\!]\cdot
[\![\,t_{j,m}\in W(t_{i,n},\tau)\,]\!],
\end{equation}
where $w_{j,m}=\upsilon_{j,m}p_{j,m}$ gives the notional feature and
$w_{j,m}=1$ the trade count, with $\upsilon$ the quantity, $p$ the price,
$\varepsilon$ the direction and $[\![\cdot]\!]$ the Iverson bracket. ``Before''
is $t_{j,m}\in(t_{i,n}-\tau,\,t_{i,n})$ and ``after'' is
$t_{j,m}\in(t_{i,n},\,t_{i,n}+\tau)$. The calibrated window is
$\tau^\star=707\,\mu\mathrm{s}$, fitted on 2017 SPY data.

\subsection{Normalisation, reduction, clustering}

\paragraph{Percentile rank against comparable history.} Each feature is replaced
by its rank within a 90-day lookback restricted to the \emph{same time-of-day
bucket} (65-minute buckets, six per session):
\begin{equation}
\tilde x^{(j)}_{i,n}=\frac{\sum_{m\in\mathcal{W}_{i,n}}[\![\,x^{(j)}_{i,m}<x^{(j)}_{i,n}\,]\!]}{|\mathcal{W}_{i,n}|-1}\in[0,1].
\end{equation}
This is the step that makes a microstructure feature comparable across days and
across times of day, and it is the one most often skipped.

\paragraph{PCA.} On mean-centred normalised features (centred, not
variance-standardised); $M$ is the smallest number of components with at least
90\% explained variance, calibrated to $M^\star=5$. Components are
sign-aligned across windows,
$v_{j,d+1}\leftarrow\sgn(v_{j,d+1}^{\top}v_{j,d})\,v_{j,d+1}$, without which
cluster labels flip arbitrarily.

\paragraph{Clustering.} Mini-batch $k$-means++ on the component scores in a
14-day sliding window advancing daily: 100 initialisations, at most 100 passes,
mini-batches of $2^{15}$ trades, lowest inertia kept, assignments retained only
for the last day. Both the Calinski--Harabasz score and the Davies--Bouldin
index select $K^\star=2$. Labels are matched across consecutive windows by
\begin{equation}
\pi^\star_{d+1}=\arg\min_{\pi\in\mathrm{Sym}(K)}\sum_{k=1}^{K}\lVert m_{\pi(k),d+1}-m_{k,d}\rVert^2 .
\end{equation}

\paragraph{Cluster-conditional order imbalance.} With $V^{\pm}$ the buy and sell
volumes,
\begin{equation}
\mathrm{OI}_{i,d}=\frac{V^{+}_{i,d}-V^{-}_{i,d}}{V^{+}_{i,d}+V^{-}_{i,d}},
\end{equation}
and $\mathrm{COI}_{i,k,d}$ the same using only trades in cluster $k$. These are
orthogonalised against a seven-factor loading matrix (Fama--French five, UMD,
and plain OI) and turned into factor-mimicking portfolios by minimum-variance
weights subject to unit exposure to the target and zero to the rest.

\paragraph{Results.} SPY and constituents, 2018--2021. At $K=2$ the first
cluster's factor-mimicking portfolio earns an annualised Sharpe ratio of
1.91, significant at 1\% after Bonferroni correction; the second earns $-0.69$.
At $K=5$ only cluster 3 is significant at 5\% (1.52). The claim to take away is
narrow: \emph{some} clusters carry information beyond raw order imbalance.

\begin{lstlisting}[caption={The Petit pipeline: percentile normalisation, PCA with sign
alignment, mini-batch k-means with label matching, cluster-conditional imbalance.}]
from sklearn.cluster import MiniBatchKMeans
from itertools import permutations


def percentile_normalise(feat, bucket_of_day, lookback_days=90):
    """Rank each feature within the same time-of-day bucket over 90 days."""
    out = pd.DataFrame(index=feat.index, columns=feat.columns, dtype=float)
    for b, idx in feat.groupby(bucket_of_day).groups.items():
        sub = feat.loc[idx].sort_index()
        win = f"{lookback_days}D"
        out.loc[idx] = sub.rolling(win, closed="left").apply(
            lambda w: (w[:-1] < w[-1]).sum() / max(len(w) - 1, 1), raw=True)
    return out


def align_components(new_v, old_v):
    """Sign-align PCs so cluster labels do not flip between windows."""
    return np.array([np.sign(n @ o) * n for n, o in zip(new_v, old_v)])


def match_labels(new_centres, old_centres):
    """Permutation of new cluster labels closest to the previous window."""
    best, best_cost = None, np.inf
    for perm in permutations(range(len(new_centres))):
        cost = sum(np.sum((new_centres[p] - old_centres[k]) ** 2)
                   for k, p in enumerate(perm))
        if cost < best_cost:
            best, best_cost = perm, cost
    return best


def cluster_conditional_oi(trades, labels, k):
    """COI for cluster k: signed volume share within that cluster."""
    sub = trades[labels == k]
    vp = sub.loc[sub["sign"] > 0, "volume"].sum()
    vm = sub.loc[sub["sign"] < 0, "volume"].sum()
    return (vp - vm) / (vp + vm) if (vp + vm) > 0 else np.nan
\end{lstlisting}

%======================================================================
\section{Prerequisites for corporate bonds}\label{sec:bonds}

Two things must be in place before any of the above is computed on a bond tape:
the trades must be signed, and the prices must be corrected for bid-ask noise.

\subsection{Signing trades}

\textbf{Source.} Fedenia, Ronen \& Nam (2021), \emph{Machine Learning in the
Corporate Bond Market: A New Classifier}, SSRN 3848068.

Every imbalance feature in this document needs $\varepsilon_t$. On TRACE the
sign is not published in the standard feed, so it has to be inferred. The
classical rules are:

\begin{itemize}[leftmargin=2em,itemsep=1pt]
\item \textbf{Tick rule.} Buy if the previous price was lower, sell if higher;
on a zero tick, compare with the nearest previous non-zero tick.
\item \textbf{Quote rule.} Buy above the quote midpoint, sell below.
\item \textbf{Lee--Ready.} Tick rule at the midpoint, quote rule outside it.
\item \textbf{Bulk volume classification.} Aggregate into bars and split volume
by $\hat V^{B}_{\tau}=V_{\tau}\,Z\bigl((P_{\tau}-P_{\tau-1})/\sigma_{\Delta P}\bigr)$
with $Z$ a Student-$t$ CDF.
\end{itemize}

The paper trains a random forest on labelled TRACE Enhanced data, where the true
initiator is known, and beats all of them: accuracy 81.1\% for the best feature
set against 68.8\% for the tick rule across roughly 143 million trades, July
2002 to December 2019. On equities the same approach reaches 77.6\% against
74.0\% for Lee--Ready. The model is deliberately plain --- 100 trees, one leaf
--- fitted on fifteen rolling three-year windows and tested on the following one
and three years, which is the time-ordered split that matters.

Feature groups that carry the accuracy: price and five lags; time of day
decomposed into hour, minute, second, weekday, month, quarter; volume and five
lags; the tick-rule indicator itself as an input; running and rolling
statistics over the last 20 trades; elapsed time since the last trade and since
the last price change; and, where available, quotes.

\paragraph{Implementation warning.} The labelled indicator exists in the
Enhanced historical file, not in the standard feed. A classifier trained on one
and applied to the other must respect that the feature set differs.

\subsection{Correcting for bid-ask noise}

\textbf{Source.} Dickerson, Robotti \& Rossetti (2024), \emph{Common Pitfalls in
the Evaluation of Corporate Bond Strategies}, SSRN 4575879.

A price-based signal and the return that follows it share a price. Bid-ask
bounce in that shared price creates predictability that is pure arithmetic. The
correction is an \emph{implementation gap}: build the signal from a price at
least one trading day before the price that starts the return.

With $p_{t,d}$ the price on the first available trading day of month $t$ and
$p_{t,d+n}$ a price at least one trading day before the month-end price $P_t$,
\begin{equation}
STREV^{*}_{i,t}=\frac{p_{i,t,d+n}+c_{i,t}}{p_{i,t,d}}-1
\qquad\text{against}\qquad
STREV^{MMN}_{i,t}=\frac{P_{i,t}+c_{i,t}}{P_{i,t-1}}-1,
\end{equation}
where $c_{i,t}$ is the coupon and all prices include accrued interest. The
realised return is
\begin{equation}
r^{B}_{i,t+1}=\frac{P^{f}_{i,t+1}+C_{i,t+1}}{P^{f}_{i,t}}-1,
\end{equation}
with $P^{f}$ the full price. The month-end price must fall in the last five
business days of the month, and extending the gap beyond one business day
changes little.

\paragraph{Effect.} The value-weighted short-term reversal premium falls from
0.90\% per month ($t=4.79$) to 0.03\%, an insignificant residue --- a reduction
over 90\%. Alphas of other price-based sorts fall by 38\% to 68\%.

\paragraph{Ex ante versus ex post.} Ex-post filtering sets the universe using a
percentile of the pooled return distribution across all bonds and months, which
uses information from after formation. Ex-ante filtering uses only information
available at or before formation month $t$. The difference is not cosmetic: the
Sharpe ratio of MOM(3,3) quintuples from 0.08 to 0.54 under the infeasible
version, and MOM(6,6) goes from $-0.06\%$ per month feasible to $0.33\%$
infeasible.

Corrected data and the construction code are public:
\url{https://openbondassetpricing.com} and the PyBondLab package.

\begin{lstlisting}[caption={Implementation gap and a signed bond tape.}]
def implementation_gap_signal(prices, coupons, min_gap_days=1):
    """Signal price must precede the return's starting price by >= 1 trading day.

    `prices` is a per-bond daily series with a DatetimeIndex.
    """
    monthly = prices.resample("ME")
    first_px = monthly.first()                     # p_{t,d}
    signal_px = prices.shift(min_gap_days).resample("ME").last()   # p_{t,d+n}
    ret_start = monthly.last()                     # P_t, starts the realised return
    signal = (signal_px + coupons.resample("ME").sum()) / first_px - 1.0
    # ret_start is returned so the caller can check that the signal price is
    # strictly earlier than the price that starts the realised return.
    return signal.rename("strev_corrected"), ret_start.rename("return_start_px")


def bond_flow_features(trades, bucket="1D"):
    """Trade-only imbalance features for a bond tape.

    Requires a signed tape (see section on signing). Uses the generalised
    imbalance of section 7, which needs no order book.
    """
    g = trades.groupby([pd.Grouper(freq=bucket), "isin"])
    vol = g["volume"].sum()
    signed = g.apply(lambda d: float((np.sign(d["sign"]) * d["volume"]).sum()))
    out = pd.DataFrame({"signed_volume": signed, "volume": vol})
    out["oi"] = out["signed_volume"] / out["volume"].replace(0.0, np.nan)
    for a in (0.0, 0.5, 1.0):
        out[f"imb_a{a}"] = g.apply(
            lambda d, a=a: float((np.sign(d["sign"]) * d["volume"] ** a).sum()))
    return out
\end{lstlisting}

%======================================================================
\section{Supporting results}\label{sec:support}

\paragraph{Long memory of order flow.} Gould, Porter \& Howison (2016) estimate
the Hurst exponent of order-sign series in FX spot at $H\approx0.7$ across
EUR/USD, GBP/USD and EUR/GBP, using detrended fluctuation analysis and
log-periodogram regression, with $H=1-\gamma/2$ linking it to the
autocorrelation decay exponent. \textbf{Their sign convention is inverted}:
$+1$ denotes a sell. They report $H$ only, not $\gamma$.

\paragraph{Cross-response between stocks.} Wang \& Guhr (2018) define
\begin{equation}
R_{ij}(\tau)=\bigl\langle r_i(t,\tau)\,\varepsilon_j(t)\bigr\rangle_t,\qquad i\neq j,
\end{equation}
with $r_i$ the log mid-price return and $\varepsilon_j\in\{-1,0,+1\}$ the trade
sign, signed from price changes rather than by Lee--Ready. There is no
normalisation by volume in the definition. The finding relevant here: self-impact
has a permanent component of order $10^{-11}$ against a temporary component of
order $10^{-4}$, and cross-impacts have \emph{no} permanent component. Volume
impact scales as $v^{0.51}$.

\paragraph{Fill probabilities.} Lokin \& Yu (2024) model limit order fill
probabilities with state-dependent intensities, conditioning on the distance to
the opposite best quote and the spread. Despite the common citation,
\textbf{this paper defines no queue-imbalance variable}: it states only that
imbalance is implicit in conditioning on both sides' queue sizes. If you need a
queue-imbalance definition, it has to come from elsewhere.

%======================================================================
\section{Source map}\label{sec:map}

\begin{table}[h]
\centering
\captionsetup{font=small}
\caption{What each paper contributes, and what it needs.}
\small
\begin{tabular}{p{4.2cm}p{5.2cm}p{4.6cm}}
\toprule
Paper & Contribution & Data required \\
\midrule
Xu, Gould \& Howison (2020) & Multi-level OFI vector; ridge over OLS & Full-depth book events \\
Cont, Cucuringu \& Zhang (2023) & Depth scaling; integrated OFI via PCA; cross-impact & Full-depth book events, many names \\
Sitaru, Calinescu \& Cucuringu (2023) & Add/cancel/trade decomposition & Market-by-order \\
Su et al.\ (2021) & Generalised OFI for snapshots; log transform & L2 snapshots \\
Hu \& Zhang (2025) & Shock/horizon framing; indicator screening test & Snapshots, one instrument \\
Maitrier \& Bouchaud (2025) + companion & Generalised imbalance $I^a$; hump in $a$; proxy metaorders & Signed trade tape only \\
Petit, Cucuringu \& Cartea (2025) & 16-feature construction; percentile normalisation; clustering & Trade tape, ETF + constituents \\
Fedenia, Ronen \& Nam (2021) & Trade signing by random forest & Labelled TRACE Enhanced \\
Dickerson et al.\ (2024) & Implementation gap; ex-ante filtering & Bond tape with accrued interest \\
Gould et al.; Wang \& Guhr; Lokin \& Yu & Long memory; cross-response; fill probabilities & Various \\
\bottomrule
\end{tabular}
\end{table}

\paragraph{What transfers to a bond tape.} Sections~\ref{sec:mlofi}
through~\ref{sec:hu} need a limit order book and do not transfer directly.
Section~\ref{sec:genimb} needs only signed trades and sizes, and
section~\ref{sec:petit} needs only a trade tape plus a grouping structure. Those
two are the portable ones, and both depend on section~\ref{sec:bonds}: sign the
trades first, and put a gap between the signal price and the return price.

\end{document}
